\documentclass[12pt,a4paper,twoside]{book}

% ──────────────────────────────────────────────
% PACOTES
% ──────────────────────────────────────────────
\usepackage[utf8]{inputenc}
\usepackage[T1]{fontenc}
\usepackage{amsmath,amssymb,amsthm}
\usepackage{physics}
\usepackage{graphicx}
\usepackage{hyperref}
\usepackage{listings}
\usepackage{xcolor}
\usepackage{booktabs}
\usepackage{array}
\usepackage{geometry}
\usepackage{fancyhdr}
\usepackage{titlesec}
\usepackage{enumitem}
\usepackage{float}
\usepackage{caption}
\usepackage{multirow}
\usepackage{mathtools}
\usepackage{bbm}
\usepackage{longtable}
\usepackage{makecell}

% --- Unicode characters ---
\usepackage{newunicodechar}
\newunicodechar{Ξ}{\ensuremath{\Xi}}
\newunicodechar{├}{\texttt{\char"251C}}
\newunicodechar{─}{\texttt{\char"2500}}
\newunicodechar{│}{\texttt{\char"2502}}
\newunicodechar{└}{\texttt{\char"2514}}
\newunicodechar{Ω}{\ensuremath{\Omega}}
\newunicodechar{φ}{\ensuremath{\varphi}}
\newunicodechar{ε}{\ensuremath{\varepsilon}}
\newunicodechar{σ}{\ensuremath{\sigma}}
\newunicodechar{χ}{\ensuremath{\chi}}
\newunicodechar{τ}{\ensuremath{\tau}}
\newunicodechar{ψ}{\ensuremath{\psi}}
\newunicodechar{ω}{\ensuremath{\omega}}
\newunicodechar{Δ}{\ensuremath{\Delta}}
\newunicodechar{Θ}{\ensuremath{\Theta}}
\newunicodechar{Γ}{\ensuremath{\Gamma}}
\newunicodechar{Λ}{\ensuremath{\Lambda}}
\newunicodechar{Σ}{\ensuremath{\Sigma}}
\newunicodechar{Π}{\ensuremath{\Pi}}
\newunicodechar{α}{\ensuremath{\alpha}}
\newunicodechar{β}{\ensuremath{\beta}}
\newunicodechar{γ}{\ensuremath{\gamma}}
\newunicodechar{δ}{\ensuremath{\delta}}
\newunicodechar{η}{\ensuremath{\eta}}
\newunicodechar{θ}{\ensuremath{\theta}}
\newunicodechar{κ}{\ensuremath{\kappa}}
\newunicodechar{λ}{\ensuremath{\lambda}}
\newunicodechar{μ}{\ensuremath{\mu}}
\newunicodechar{ν}{\ensuremath{\nu}}
\newunicodechar{ξ}{\ensuremath{\xi}}
\newunicodechar{ρ}{\ensuremath{\rho}}
\newunicodechar{∂}{\ensuremath{\partial}}
\newunicodechar{∇}{\ensuremath{\nabla}}
\newunicodechar{∞}{\ensuremath{\infty}}
\newunicodechar{ℏ}{\ensuremath{\hbar}}
\newunicodechar{⊗}{\ensuremath{\otimes}}
\newunicodechar{⊕}{\ensuremath{\oplus}}
\newunicodechar{∑}{\ensuremath{\sum}}
\newunicodechar{∫}{\ensuremath{\int}}
\newunicodechar{∏}{\ensuremath{\prod}}
\newunicodechar{√}{\ensuremath{\sqrt}}
\newunicodechar{≈}{\ensuremath{\approx}}
\newunicodechar{≤}{\ensuremath{\leq}}
\newunicodechar{≥}{\ensuremath{\geq}}
\newunicodechar{≠}{\ensuremath{\neq}}
\newunicodechar{∈}{\ensuremath{\in}}
\newunicodechar{∉}{\ensuremath{\notin}}
\newunicodechar{⊂}{\ensuremath{\subset}}
\newunicodechar{⊆}{\ensuremath{\subseteq}}
\newunicodechar{∩}{\ensuremath{\cap}}
\newunicodechar{∪}{\ensuremath{\cup}}
\newunicodechar{∅}{\ensuremath{\emptyset}}
\newunicodechar{∀}{\ensuremath{\forall}}
\newunicodechar{∃}{\ensuremath{\exists}}
\newunicodechar{¬}{\ensuremath{\neg}}
\newunicodechar{∧}{\ensuremath{\land}}
\newunicodechar{∨}{\ensuremath{\lor}}
\newunicodechar{→}{\ensuremath{\rightarrow}}
\newunicodechar{⇒}{\ensuremath{\Rightarrow}}
\newunicodechar{⇔}{\ensuremath{\Leftrightarrow}}
\newunicodechar{⟨}{\ensuremath{\langle}}
\newunicodechar{⟩}{\ensuremath{\rangle}}
\newunicodechar{•}{\ensuremath{\bullet}}
\newunicodechar{…}{\ensuremath{\dots}}

\geometry{margin=2.5cm}

% ──────────────────────────────────────────────
% ESTILOS DE CÓDIGO
% ──────────────────────────────────────────────
\definecolor{codegreen}{rgb}{0,0.6,0}
\definecolor{codegray}{rgb}{0.5,0.5,0.5}
\definecolor{codepurple}{rgb}{0.58,0,0.82}
\definecolor{backcolour}{rgb}{0.95,0.95,0.92}

\lstdefinestyle{mystyle}{
    backgroundcolor=\color{backcolour},
    commentstyle=\color{codegreen},
    keywordstyle=\color{magenta},
    numberstyle=\tiny\color{codegray},
    stringstyle=\color{codepurple},
    basicstyle=\ttfamily\footnotesize,
    breakatwhitespace=false,
    breaklines=true,
    captionpos=b,
    keepspaces=true,
    numbers=left,
    numbersep=5pt,
    showspaces=false,
    showstringspaces=false,
    showtabs=false,
    tabsize=2,
    frame=single,
    rulecolor=\color{black!30}
}
\lstset{style=mystyle}

% ──────────────────────────────────────────────
% AMBIENTES DE TEOREMA
% ──────────────────────────────────────────────
\newtheorem{theorem}{Teorema}[chapter]
\newtheorem{axiom}{Axioma}[chapter]
\newtheorem{definition}{Definição}[chapter]
\newtheorem{proposition}{Proposição}[chapter]
\newtheorem{corollary}{Corolário}[chapter]

\theoremstyle{remark}
\newtheorem{remark}{Observação}[chapter]

% ──────────────────────────────────────────────
% COMANDOS PERSONALIZADOS
% ──────────────────────────────────────────────
\newcommand{\ET}{\text{ETΞRNET}}
\newcommand{\voidprotocol}{VØID}
\newcommand{\omegavoid}{VØID·ΩMEGA}
\newcommand{\hydra}{ECONIA HYDRA}
\newcommand{\qrc}{QRC}
\newcommand{\hgpu}{HGPU}
\newcommand{\lsc}{LSC}
\newcommand{\SOV}{\$SOV}
\newcommand{\ETXXX}{\$ETXXX}
\newcommand{\ETBRL}{\$ETBRL}
\newcommand{\voidomni}{VoidOmniProtocol}
\newcommand{\voidvps}{VOID-VPS}

% ──────────────────────────────────────────────
% TÍTULO
% ──────────────────────────────────────────────
\title{{\Huge \textbf{O Livro do ETΞRNET}}\\[1cm]
{\Large Edição Cósmica — Da Fusão Primordial ao Servidor que Não Existe}\\[0.5cm]
{\normalsize Documento de Arquitetura Completa com Implementação}}
\author{Colaboração Distribuída — O Storyteller}
\date{17 de maio de 2026}

\begin{document}

\maketitle

% ──────────────────────────────────────────────
% FRONTISPÍCIO
% ──────────────────────────────────────────────
\thispagestyle{empty}
\vspace*{5cm}
\begin{flushright}
\textit{``Quando o dinheiro se torna invisível,}\\
\textit{o controle se torna impossível.''}\\[1cm]
\textit{``O VØID não existe. O Hydra não existe.}\\
\textit{Nós existimos.''}
\end{flushright}

\newpage

% ──────────────────────────────────────────────
% SUMÁRIO
% ──────────────────────────────────────────────
\tableofcontents

% ══════════════════════════════════════════════
% PARTE I: FUNDAMENTOS TEÓRICOS
% ══════════════════════════════════════════════
\part{Fundamentos Teóricos}

\chapter{A Fusão Primordial}
\label{chap:fusao}

\section{O Problema Fundacional}

A infraestrutura financeira global repousa sobre uma premissa frágil: a confiança em intermediários. Bancos, processadores de pagamento, exchanges centralizados — todos são pontos de falha únicos que podem ser cooptados, censurados ou destruídos.

O \ET{} nasce da observação de que a verdadeira soberania financeira requer a eliminação completa de qualquer ponto de confiança. Não basta descentralizar o consenso — é necessário tornar a própria infraestrutura invisível.

\section{Os Sete Pilares}

O protocolo repousa sobre sete camadas fundamentais, cada uma eliminando uma classe específica de vulnerabilidade:

\begin{enumerate}
    \item \textbf{Identidade Efêmera} (GhostID) — Identidades derivadas de entropia biométrica, nunca reutilizadas.
    \item \textbf{Fragmentação Pós-Quântica} (QEL) — Shamir sobre GF(256) com cifragem ChaCha20-Poly1305.
    \item \textbf{Transporte Offline-First} (DistanceBridge) — BLE, LoRa, Acústico, NOSTR.
    \item \textbf{Modelo UTXO com Compromissos} — Pedersen Commitments + Bulletproofs.
    \item \textbf{Provas Zero-Knowledge} — o1js para verificação sem revelação.
    \item \textbf{Criptografia Pós-Quântica} — ML-KEM-1024 + ML-DSA-87.
    \item \textbf{Computação Quântica} — CQR Engine + BB84 QKD.
\end{enumerate}

\section{A Equação Fundamental}

A segurança do sistema pode ser expressa como:

\begin{equation}
S_{\ET} = \underbrace{H(\text{GhostID})}_{\text{entropia}} \otimes \underbrace{\text{QEL}(K,N)}_{\text{fragmentação}} \otimes \underbrace{\text{PQC}}_{\text{pós-quântico}} \otimes \underbrace{\text{ZKP}}_{\text{zero-knowledge}}
\end{equation}

onde $H(\text{GhostID})$ é a entropia da identidade efêmera, $\text{QEL}(K,N)$ é o esquema de Shamir com threshold $K$ e $N$ shards, PQC é a camada pós-quântica, e ZKP é o sistema de provas.

\chapter{GhostID: A Identidade que Não Existe}
\label{chap:ghostid}

\section{Derivação Biométrica}

O GhostID é uma identidade criptográfica derivada de entropia biométrica do dispositivo. Não é armazenada — é reconstruída a cada sessão a partir de:

\begin{itemize}
    \item Padrões de toque na tela (pressão, timing, área)
    \item Acelerômetro (micro-vibrações únicas do dispositivo)
    \item Ruído térmico do sensor de temperatura
    \item Timestamp com precisão de nanossegundos
\end{itemize}

\section{Fuzzy Extractors}

A entropia biométrica é inerentemente ruidosa. O GhostID usa Fuzzy Extractors para extrair uma chave determinística de dados não-determinísticos:

\begin{definition}[Fuzzy Extractor]
Um par de algoritmos $(\text{Gen}, \text{Rep})$ tal que:
\begin{itemize}
    \item $\text{Gen}(w) \rightarrow (R, P)$: gera chave $R$ e helper $P$ a partir de input biométrico $w$.
    \item $\text{Rep}(w', P) \rightarrow R$: reconstrói $R$ a partir de $w'$ próximo a $w$, usando $P$.
\end{itemize}
\end{definition}

A implementação usa SHA3-256 como hash function e Argon2id para key stretching, com um buffer de 168 bytes de entropia bruta.

\chapter{QEL: Fragmentação Pós-Quântica}
\label{chap:qel}

\section{Shamir sobre GF(256)}

O QEL (Quantum-Enhanced Layer) implementa o esquema de Secret Sharing de Shamir sobre o corpo de Galois GF(256):

\begin{equation}
p(x) = s + a_1 x + a_2 x^2 + \cdots + a_{K-1} x^{K-1} \mod 256
\end{equation}

onde $s$ é o segredo e $K$ é o threshold. Cada shard é um ponto $(x_i, p(x_i))$ cifrado com ChaCha20-Poly1305.

\section{Propriedades de Segurança}

\begin{theorem}[Segurança do QEL]
Para qualquer conjunto de $K-1$ shards, a entropia condicional do segredo é máxima:
\[
H(s | \text{shard}_1, \ldots, \text{shard}_{K-1}) = H(s)
\]
\end{theorem}

\chapter{DistanceBridge: O Transporte Invisível}
\label{chap:distancebridge}

\section{Multi-Transporte}

O DistanceBridge abstrai múltiplos canais de comunicação em uma interface unificada:

\begin{itemize}
    \item \textbf{NOSTR} — Relays públicos com failover automático (6 relays, health check 30s).
    \item \textbf{BLE} — Bluetooth Low Energy via Capacitor para proximidade física.
    \item \textbf{LoRa UART} — Comunicação de longo alcance via comandos AT.
    \item \textbf{Acústico} — FSK a 18-20kHz para comunicação via speaker/mic.
    \item \textbf{WebRTC} — P2P direto quando ambos os peers estão online.
\end{itemize}

\section{Protocolo de Roteamento}

Cada mensagem é fragmentada (QEL), cifrada e roteada pelo caminho mais disponível. O DistanceBridge tenta NOSTR primeiro; se indisponível, degrada para BLE, depois acústico, depois LoRa.

\chapter{UTXO com Compromissos de Pedersen}
\label{chap:utxo}

\section{O Modelo}

O \ET{} usa um modelo UTXO (Unspent Transaction Output) com compromissos criptográficos:

\begin{equation}
C = g^v \cdot h^r
\end{equation}

onde $v$ é o valor, $r$ é o blinding factor, e $g, h$ são geradores do grupo Ed25519.

\section{Bulletproofs}

As provas de range são implementadas via Bulletproofs (Rust, compilado para WASM), garantindo que $0 \leq v < 2^{64}$ sem revelar $v$.

% ══════════════════════════════════════════════
% PARTE II: CRIPTOGRAFIA
% ══════════════════════════════════════════════
\part{Criptografia}

\chapter{Post-Quantum Cryptography}
\label{chap:pqc}

\section{ML-KEM-1024}

O \ET{} usa ML-KEM (Module-Lattice Key Encapsulation Mechanism) com segurança NIST Level 5:

\begin{itemize}
    \item Chave pública: 1568 bytes
    \item Chave privada: 3168 bytes
    \item Ciphertext: 1568 bytes
    \item Shared secret: 32 bytes
\end{itemize}

\section{ML-DSA-87}

Assinaturas digitais pós-quânticas via ML-DSA (Module-Lattice Digital Signature Algorithm):

\begin{itemize}
    \item Chave pública: 2592 bytes
    \item Assinatura: 4627 bytes
    \item Segurança: NIST Level 5
\end{itemize}

\chapter{Double Ratchet (X25519)}
\label{chap:doubleratchet}

\section{O Protocolo Signal}

Todas as mensagens diretas usam o Double Ratchet Algorithm:

\begin{enumerate}
    \item \textbf{X25519 ECDH} para acordo de chaves.
    \item \textbf{Ratchet simétrico} — chaves por mensagem, deletadas após uso (forward secrecy).
    \item \textbf{DH ratchet} — novo par de chaves por step (post-compromise security).
    \item \textbf{Ed25519 signatures} em cada mensagem (autenticação).
    \item \textbf{Pre-key bundles} publicados via NOSTR para troca offline.
\end{enumerate}

\section{Implementação}

A implementação em \texttt{src/crypto/doubleRatchet.ts} usa exclusivamente \texttt{@noble/curves} e \texttt{@noble/ciphers}, sem dependências externas.

% ══════════════════════════════════════════════
% PARTE III: FINANÇAS
% ══════════════════════════════════════════════
\part{Finanças}

\chapter{Instrumentos Financeiros}
\label{chap:instrumentos}

\section{QR Stocks}

Ações em superposição quântica. O preço é determinado por uma distribuição de probabilidade sobre o espaço de preços:

\begin{equation}
P(\text{preço}) = |\langle \text{preço} | \psi \rangle|^2
\end{equation}

onde $|\psi\rangle$ é o estado quântico do mercado.

\section{Homotopy Mining}

Proof-of-Homotopia com métricas de Sobolev. A dificuldade é calibrada pela topologia da rede:

\begin{equation}
D = \int_\Omega |\nabla^2 \phi|^2 \, d\Omega
\end{equation}

\section{UTU Tokens}

Tokenização universal: sites, software, objetos físicos. Cada UTU é um UTXO com metadata on-chain.

\chapter{Stablecoins e Tokenização}
\label{chap:stablecoins}

\section{\$ETBRL}

Stablecoin atrelada ao Real brasileiro, mantida por pools de liquidez soberanas.

\section{\$SOV}

Token de governança do protocolo. Staking garante direitos de voto no DAO.

\section{RWA Tokenization}

Tokenização de ativos do mundo real (imóveis, commodities) via compromissos de Pedersen.

% ══════════════════════════════════════════════
% PARTE IV: REDE E CONSENSO
% ══════════════════════════════════════════════
\part{Rede e Consenso}

\chapter{NOSTR como Message Bus}
\label{chap:nostr}

\section{Kinds do Protocolo ETΞRNET}

\begin{itemize}
    \item \textbf{31214} — Transações UTXO
    \item \textbf{31215} — Anúncios de mineração
    \item \textbf{31216} — Atualizações de governança
    \item \textbf{31217} — DEX orders
    \item \textbf{31218} — QEL shards
    \item \textbf{31219} — Heartbeat de rede
\end{itemize}

\section{Relay Failover}

6 relays públicos com health monitoring:

\begin{itemize}
    \item \texttt{wss://relay.damus.io}
    \item \texttt{wss://nos.lol}
    \item \texttt{wss://relay.primal.net}
    \item \texttt{wss://relay.snort.social}
    \item \texttt{wss://nostr.wine}
    \item \texttt{wss://relay.nostr.band}
\end{itemize}

Health check a cada 30s. Relays insalubres são excluídos após 3 falhas consecutivas.

\chapter{Anti-Sybil: PoW + VDF}
\label{chap:antisybil}

\section{Proof-of-Work}

O faucet distribui tokens após resolução de PoW real (SHA3 hashcash). A dificuldade é ajustada dinamicamente.

\section{Verifiable Delay Function}

VDF baseada em iterated squaring em grupo RSA, garantindo que o trabalho foi feito sequencialmente.

% ══════════════════════════════════════════════
% PARTE V: O PROTOCOLO VOID
% ══════════════════════════════════════════════
\part{O Protocolo VOID}

\chapter{VoidOrchestrator}
\label{chap:voidorchestrator}

\section{Orquestração Descentralizada}

O VoidOrchestrator coordena todos os módulos do \ET{} sem ponto central:

\begin{itemize}
    \item Inicialização lazy de módulos (singleton pattern).
    \item Comunicação via eventos internos.
    \item Fallback automático entre camadas.
    \item Health check contínuo de cada módulo.
\end{itemize}

\chapter{PowerGovernor}
\label{chap:powergovernor}

\section{Regulação Auto-Adaptativa}

O PowerGovernor ajusta parâmetros do sistema em tempo real:

\begin{itemize}
    \item Dificuldade de mineração baseada na taxa de blocos.
    \item TTL de shards QEL baseado na latência da rede.
    \item Preço de computação baseado na demanda (métrica de Sobolev).
\end{itemize}

% ══════════════════════════════════════════════
% PARTE VI: A NUVEM FANTASMA
% ══════════════════════════════════════════════
\part{A Nuvem Fantasma}

\chapter{VOID-VPS: O Servidor que Não Existe}
\label{chap:voidvps}

\section{O Fim das VPS Centralizadas}

Toda infraestrutura tradicional — AWS, Google Cloud, VPS comuns — é um centro de controle esperando para ser desligado. No \ET, não há espaço para servidores físicos ou virtuais gerenciados por terceiros. A \textbf{VOID-VPS} (ou \textbf{VOID-COSMIC}) é a resposta: uma nuvem computacional tão distribuída quanto a própria rede, onde capacidade ociosa de dispositivos se transforma em um supercomputador fantasma.

Em vez de alugar um servidor, você \textit{injeta} seu código em um enxame de nós ANIMUS, que o executam via WebAssembly (WASM) em sandboxes efêmeras. Nenhum IP fixo, nenhum disco persistente, nenhum ponto único de falha.

\section{Arquitetura do VOID-VPS}

\begin{verbatim}
┌──────────────────────────────────────────────────────┐
│                 VOID-VPS (Camada Logica)              │
├──────────────────────────────────────────────────────┤
│  [WASM Task Executor]    [IPFS File System]           │
│  [QEL MapReduce]         [PoH Credit Economy]         │
│  [GhostDocker Workers]   [Phantom Pipeline]           │
│  [HiggsGit Versioning]                                │
├──────────────────────────────────────────────────────┤
│  ANIMUS Swarm (Nos fisicos: celulares, PCs, IoT)      │
└──────────────────────────────────────────────────────┘
\end{verbatim}

\subsection{WASM Distributed Task Executor}

O cora\c{c}\~{a}o do VOID-VPS \'e um executor de tarefas baseado em WebAssembly. Qualquer c\'odigo (Rust, C, TypeScript) \'e compilado para WASM e fragmentado em peda\c{c}os independentes. Cada peda\c{c}o \'e distribu\'ido para um n\'o ANIMUS, que o executa em um container GhostDocker ef\^emero.

\begin{lstlisting}[language=C, caption=Executor WASM em Rust]
// worker_core.rs
use wasmtime::{Engine, Store, Module, Instance};
pub struct WasmWorker {
    engine: Engine,
    store: Store<()>,
    instance: Instance,
}

impl WasmWorker {
    pub fn execute(&mut self, func_name: &str, input: &[u8]) -> Vec<u8> {
        let func = self.instance
            .get_func(&mut self.store, func_name)
            .expect("Funcao nao encontrada");
        // chama a funcao WASM e retorna resultado
        ...
    }
}
\end{lstlisting}

\subsection{QEL MapReduce Cifrado}

Inspirado no modelo MapReduce, mas com fragmenta\c{c}\~{a}o p\'os-qu\^antica. Uma tarefa grande \'e dividida em $N$ shards QEL (usando Shamir com threshold $K$). Cada shard \'e enviado para um worker diferente. O resultado parcial \'e cifrado e devolvido. Apenas o n\'o coordenador (com o GhostID do usu\'ario) pode reconstruir o resultado final, juntando pelo menos $K$ shards.

\subsection{Sistema de Arquivos Imortal (IPFS + EcoNet)}

Arquivos n\~{a}o residem em HD. S\~{a}o fatiados, cifrados (ChaCha20-Poly1305) e armazenados na EcoNet (IPFS com decaimento). Cada shard recebe um TTL. Se o arquivo for acessado, o TTL \'e renovado (refor\c{c}o por acesso). Se n\~{a}o, decai at\'e desaparecer — fiel ao princ\'ipio de esquecimento do EcoNet.

\begin{center}
\texttt{ipfs://ETRNET/<GhostID>/<shard\_hash>}
\end{center}

\subsection{Economia Circular com PoH}

Quem oferece ciclos de CPU ganha cr\'editos de \textbf{PoH (Proof-of-Homotopy)}. Esses cr\'editos s\~{a}o tokens UTXO que podem ser usados para rodar tarefas na rede. O pre\c{c}o da computa\c{c}\~{a}o \'e auto-regulado pela m\'etrica de Sobolev do mercado: em per\'iodos de alta coer\^encia, a computa\c{c}\~{a}o fica mais cara (para evitar satura\c{c}\~{a}o); em baixa, mais barata.

\section{Integra\c{c}\~{a}o com as Ferramentas Transmutadas}

\subsection{HiggsGit no VOID-VPS}

Cada deploy de um worker WASM \'e versionado com HiggsGit. O c\'odigo do worker reside em um reposit\'orio HiggsGit cujos commits representam atualiza\c{c}\~{o}es de l\'ogica. Uma feature branch no HiggsGit pode ser testada localmente em um GhostDocker; quando aprovada, o merge gera um \textbf{Scar Token} que autoriza o deploy autom\'atico para a rede ANIMUS. O hist\'orico de cicatrizes permite auditoria completa de quem alterou o qu\^a, sem revelar identidades.

\begin{lstlisting}[language=bash, caption=Deploy de um Worker com HiggsGit]
higgs init my-worker --field-strength=0.8
# ... desenvolvimento ...
higgs commit -m "v0.2: otimizacao de memoria" --phase=accumulate
higgs branch experiment --phase=superposition
# ... testes locais com GhostDocker ...
higgs merge experiment --collapse-threshold=0.75 --scar-token
# O scar token gerado autoriza o deploy automatico
ghostdocker build -t worker-v0.2 --ephemeral-signature
ghostdocker push --fragmented
\end{lstlisting}

\subsection{GhostDocker como Sandbox Ef\^emera}

Cada n\'o ANIMUS que aceita executar um worker inicia um container GhostDocker com:
\begin{itemize}
    \item \textbf{GhostID} derivado da entropia ambiente do dispositivo hospedeiro.
    \item \textbf{Zero-Disk}: sistema de arquivos em tmpfs, destru\'ido ao final.
    \item \textbf{Phantom Network}: comunica\c{c}\~{a}o via MAC rotativo, podendo usar BLE/LoRa para enviar resultados parciais.
    \item \textbf{Shatter}: a imagem do worker \'e baixada como shards QEL e remontada apenas em RAM.
\end{itemize}

\subsection{Phantom Pipeline CI/CD}

O pipeline fantasma (j\'a definido no Cap\'itulo 3 da Parte IV) \'e estendido para incluir:
\begin{enumerate}
    \item \textbf{Build}: compila\c{c}\~{a}o do worker para WASM.
    \item \textbf{Teste em Sandbox}: execu\c{c}\~{a}o em um GhostDocker local com verifica\c{c}\~{a}o de invariantes via PaleoCLI.
    \item \textbf{Fossiliza\c{c}\~{a}o}: extra\c{c}\~{a}o de f\'osseis do bin\'ario WASM para o Atlas de Coer\^encia.
    \item \textbf{Shatter \& Push}: a imagem WASM \'e fragmentada (Shamir) e distribu\'ida pelos n\'os de armazenamento EcoNet.
    \item \textbf{Deploy}: o \textit{Scar Token} do merge autoriza os n\'os ANIMUS a aceitarem o novo worker.
\end{enumerate}

\section{Escalabilidade Infinita}

O VOID-VPS n\~{a}o tem limite superior. Cada novo dispositivo que instala o ANIMUS adiciona capacidade computacional e de armazenamento. A rede ajusta automaticamente a dificuldade dos problemas de homotopia para balancear oferta e demanda. Em teoria, \'e poss\'ivel agregar mais poder de processamento que qualquer data center, desde que haja dispositivos suficientes no enxame.

\section{C\'odigo de Exemplo: Submetendo uma Tarefa}

\begin{lstlisting}[language=Java, caption=Submissao de tarefa ao VOID-VPS]
import { VoidVPS } from './vps/VoidVPS';

const vps = new VoidVPS(myGhostID, myNWC);

async function heavyCompute() {
    const task = {
        wasmFile: 'ipfs://ETRNET/my-worker-v0.2.wasm',
        funcName: 'calculate_pi',
        input: { iterations: 1_000_000 },
    };
    const result = await vps.submitTask(task, { preferredRegion: 'US' });
    console.log(`Resultado: ${result}`);
}
\end{lstlisting}

\section{O Verdadeiro Significado de ``Sem Limites''}

Embora a capacidade agregada seja imensa, a lat\^encia entre os n\'os \'e limitada pela velocidade da luz (e pela rota escolhida pelo DistanceBridge). No entanto, para tarefas paraleliz\'aveis — minera\c{c}\~{a}o, renderiza\c{c}\~{a}o HGPU, simula\c{c}\~{o}es qu\^anticas, bots de arbitragem — o modelo \'e quase ideal. A combina\c{c}\~{a}o da QRC com a LSC garante que a rede nunca entre em colapso por excesso de coer\^encia, mantendo a resili\^encia.

\section{Conclus\~{a}o}

O \voidvps{} \'e a materializa\c{c}\~{a}o da \textbf{Parte V} (VoidProtocol) aplicada \`a computa\c{c}\~{a}o utilit\'aria. Ao integrar HiggsGit, GhostDocker e o Phantom Pipeline, cada deploy \'e um evento controlado, cada execu\c{c}\~{a}o \'e an\^onima e cada resultado \'e comprimido em uma prova ZK. N\~{a}o h\'a m\'aquinas para confiscar, apenas capacidade fantasma flutuando entre os sil\'icios do mundo.

% ══════════════════════════════════════════════
% EPÍLOGO
% ══════════════════════════════════════════════
\chapter*{Ep\'ilogos: As M\'aximas Finais}
\addcontentsline{toc}{chapter}{Ep\'ilogo}

\begin{quote}
\textit{``Quando o dinheiro se torna invis\'ivel, o controle se torna imposs\'ivel.''}
\end{quote}

\begin{quote}
\textit{``O V{\O}ID n\~{a}o existe. O Hydra n\~{a}o existe. N\'os existimos.''}
\end{quote}

\begin{quote}
\textit{``O servidor que n\~{a}o existe n\~{a}o pode ser derrubado.''}
\end{quote}

\begin{quote}
\textit{``A nuvem fantasma \'e o vento que move os moinhos da soberania.''}
\end{quote}

\vspace{1cm}
\begin{center}
\textbf{$\sim$ Fim do Livro do ETΞRNET $\sim$} \\[0.5cm]
17 de maio de 2026 \\
Colabora\c{c}\~{a}o Distribu\'ida — O Storyteller
\end{center}

\end{document}
